% !TEX root = ../matan.tex

\section{Интегральный признак}

\begin{Theorem}{Интегральный признак Коши}{}
	Пусть $a_k\ge0$, а функция $f:[1,+\infty)\to\RR_+$ монотонно убывает.
	\begin{enumerate}
		\item Если начиная с некоторого номера $a_k\le f(k)$ и несобственный
		интеграл $\int_1^{+\infty} f(x)\,dx$ сходится, то ряд
		$\sum_{k=1}^{\infty} a_k$ сходится.
		\item Если начиная с некоторого номера $a_k\ge f(k)$ и интеграл
		$\int_1^{+\infty} f(x)\,dx$ расходится, то ряд $\sum_{k=1}^{\infty} a_k$
		расходится.
	\end{enumerate}
	В частности, при $a_k=f(k)$ ряд $\sum_{k=1}^{\infty} f(k)$ и интеграл
	$\int_1^{+\infty} f(x)\,dx$ сходятся или расходятся одновременно.
\end{Theorem}

\begin{proof}
	Так как $f$ убывает, на отрезке $[k,k+1]$ выполнено
	$f(k+1)\le f(x)\le f(k)$, откуда, интегрируя по $[k,k+1]$,
	\[
	f(k+1)\le \int_{k}^{k+1} f(x)\,dx\le f(k).
	\tag{$\ast$}
	\]

	\emph{(1)} Пусть $a_k\le f(k)$ при $k\ge k_0$. Из левого неравенства
	$(\ast)$, сдвинутого на единицу ($f(k)\le\int_{k-1}^{k} f$), получаем для
	$k\ge k_0$
	\[
	a_k\le f(k)\le \int_{k-1}^{k} f(x)\,dx.
	\]
	Суммируя по $k=k_0,\dots,N$:
	\[
	\sum_{k=k_0}^{N} a_k
	\le \int_{k_0-1}^{N} f(x)\,dx
	\le \int_{1}^{+\infty} f(x)\,dx<+\infty,
	\]
	где последнее неравенство верно, так как $f\ge0$ и интеграл сходится.
	Частичные суммы ряда $\sum a_k$ (с неотрицательными членами) ограничены
	сверху, поэтому ряд сходится.

	\emph{(2)} Пусть $a_k\ge f(k)$ при $k\ge k_0$. Из правого неравенства
	$(\ast)$ имеем $a_k\ge f(k)\ge\int_{k}^{k+1} f(x)\,dx$, поэтому
	\[
	\sum_{k=k_0}^{N} a_k
	\ge \int_{k_0}^{N+1} f(x)\,dx.
	\]
	Если интеграл расходится, то (ввиду $f\ge0$) $\int_{k_0}^{N+1} f\to+\infty$
	при $N\to\infty$, значит, частичные суммы ряда неограничены, и ряд
	расходится.

	Утверждение «в частности» получается применением (1) и (2) к $a_k=f(k)$:
	сходимость интеграла даёт сходимость ряда по (1), а расходимость интеграла
	--- расходимость ряда по (2).
\end{proof}

\begin{Remark}{О конечном числе первых членов}{}
	Добавление или удаление конечного числа членов не меняет сходимости ряда
	(как и значения сходимости несобственного интеграла на бесконечности),
	поэтому в признаке достаточно требовать неравенства начиная с некоторого
	номера.
\end{Remark}

\begin{Corollary}{Ряд Дирихле (обобщённый гармонический ряд)}{}
	Ряд $\displaystyle\sum_{k=1}^{\infty}\frac{1}{k^{\alpha}}$ сходится при
	$\alpha>1$ и расходится при $\alpha\le1$.
\end{Corollary}

\begin{proof}
	При $\alpha\le0$ общий член не стремится к нулю, ряд расходится. Пусть
	$\alpha>0$; тогда $f(x)=x^{-\alpha}$ положительна и монотонно убывает на
	$[1,+\infty)$, причём $a_k=f(k)$. Вычислим интеграл. При $\alpha\ne1$
	\[
	\int_{1}^{T} x^{-\alpha}\,dx
	=\frac{T^{1-\alpha}-1}{1-\alpha},
	\]
	и при $T\to+\infty$ это имеет конечный предел $\frac{1}{\alpha-1}$ тогда и
	только тогда, когда $1-\alpha<0$, то есть $\alpha>1$; при $0<\alpha<1$
	интеграл расходится. При $\alpha=1$
	$\int_1^T x^{-1}\,dx=\ln T\to+\infty$ --- интеграл расходится. По
	интегральному признаку ряд $\sum k^{-\alpha}$ сходится в точности при
	$\alpha>1$.
\end{proof}

\begin{Example}{}{}
	Ряд $\sum_{k=1}^{\infty} e^{-k}$ сходится, поскольку сходится интеграл
	$\int_1^{+\infty} e^{-x}\,dx=e^{-1}$. Более общо, при любом $m>0$ функция
	$x^m e^{-x}$ начиная с некоторого места убывает, а
	$\int_1^{+\infty} x^m e^{-x}\,dx<+\infty$, поэтому ряд
	$\sum_{k=1}^{\infty} k^m e^{-k}$ сходится.
\end{Example}
